\clearpage
\refstepcounter{section}
\addcontentsline{toc}{section}{\protect\numberline{}Formulario e anteprima}
\markboth{Formulario}{Formulario}
\begin{tikzpicture}
  \fill[navy,rounded corners=8pt] (0,0) rectangle (\textwidth,1.9cm);
  \node[anchor=west,font=\headfont\bfseries\fontsize{19}{23}\selectfont,text=white] at (.6cm,1.15cm) {Formulario riassuntivo};
  \node[anchor=west,font=\headlight\small,text=white!80!teal] at (.6cm,.45cm) {Le formule da sapere, in una pagina};
\end{tikzpicture}
\vspace{6pt}

\renewcommand{\arraystretch}{1.55}\small
\rowcolors{1}{white}{chiaro}
\begin{tabularx}{\linewidth}{>{\headfont\bfseries\small\color{navy}}m{3.5cm} >{\centering\arraybackslash}X}
Probabilità condizionata & $P(A|B)=\dfrac{P(A\cap B)}{P(B)}$ \qquad indipendenza: $P(A\cap B)=P(A)P(B)$\\
Probabilità totale & $p(\x)=\sum_{i=1}^{s}p(\x|w_i)P(w_i)$\\
Teorema di Bayes & $P(w_i|\x)=\dfrac{p(\x|w_i)\,P(w_i)}{p(\x)}$\\
Regola MAP & $b=\argmax_i P(w_i|\x)=\argmax_i\, p(\x|w_i)P(w_i)$\\
Probabilità di errore & $\displaystyle P(\mathrm{err})=\int_{\Re_1}p(x|w_2)P(w_2)\,dx+\int_{\Re_2}p(x|w_1)P(w_1)\,dx$\\
Stima a conteggio & $p(\x|w_i)\approx\dfrac{k_i}{n_i V}$ \qquad $P(w_i|\x)\approx\dfrac{k_i}{k}$\\
Normale 1D & $p(x)=\dfrac{1}{\sigma\sqrt{2\pi}}\,e^{-\frac{(x-\mu)^2}{2\sigma^2}}$\\
MLE 1D & $\hat\mu=\frac1n\sum_i x_i$ \qquad $\hat\sigma^2=\frac1n\sum_i(x_i-\hat\mu)^2$\\
Multinormale & $p(\x)=\dfrac{1}{(2\pi)^{d/2}|\SSigma|^{1/2}}\exp\!\big[-\tfrac12(\x-\mmu)\T\SSigma^{-1}(\x-\mmu)\big]$\\
MLE multivariata & $\mmu=\frac1n\sum_i\x_i$ \qquad $\SSigma=\frac1n\sum_i(\x_i-\mmu)(\x_i-\mmu)\T$\\
Mahalanobis & $r^2=(\x-\mmu)\T\SSigma^{-1}(\x-\mmu)$\\
Inversa $2\times2$ & $\left(\begin{smallmatrix}a&b\\c&d\end{smallmatrix}\right)^{-1}=\dfrac{1}{ad-bc}\left(\begin{smallmatrix}d&-b\\-c&a\end{smallmatrix}\right)$\\
Autovalori & $A\vv=\lambda\vv$ \qquad $\det(A-\lambda I)=0$ \qquad $\det A=\prod_k\lambda_k$\\
Discriminante & $g_i(\x)=-\frac12 r_i^2-\frac12\ln|\SSigma_i|+\ln P(w_i)$ \;(+ cost.)\\
N. parametri & $\mmu$: $d$ \qquad $\SSigma$: $\frac{d(d+1)}{2}$ \qquad Naive Bayes: $2d$\\
\end{tabularx}
\renewcommand{\arraystretch}{1.25}

\vspace{10pt}
\begin{tcolorbox}[base,colback=ambra!7,colbacktitle=ambra,coltitle=white,borderline west={3pt}{0pt}{ambra},title={Anteprima: approccio non parametrico (Parzen Window)},top=13pt]
Il sommario della lezione (slide 2 e 24) annuncia anche il metodo \textbf{Parzen Window}, che non è trattato in queste 62 slide. L'idea parte dalla formula della sezione 4, $p(\x)\approx k/(nV)$: invece di contare i punti in un cerchio fisso, si centra una piccola «finestra» (una funzione kernel, spesso gaussiana, di larghezza $h$) su \emph{ogni} campione del training set e si sommano i contributi:
\[
p(\x)\approx\frac1n\sum_{i=1}^{n}\frac{1}{h^d}\,\varphi\!\left(\frac{\x-\x_i}{h}\right)
\]
La larghezza $h$ regola il compromesso visto nella sezione 5: $h$ piccolo dà una stima piena di picchi (overfitting), $h$ grande una stima troppo liscia (underfitting).
\end{tcolorbox}

\vspace{4pt}
{\headfont\bfseries\color{navy}Per approfondire}
\begin{itemize}
  \item R. O. Duda, P. E. Hart, D. G. Stork, \emph{Pattern Classification}, 2ª ed., Wiley, 2001: capitolo 2 (teoria della decisione bayesiana) e capitolo 3 (stima MLE).
  \item K. Fukunaga, \emph{Introduction to Statistical Pattern Recognition}, Academic Press, 1990: test di normalità.
  \item S. Russell, P. Norvig, \emph{Artificial Intelligence: A Modern Approach}, 4ª ed., Pearson, 2021: capitolo sull'incertezza e gli assiomi della probabilità.
\end{itemize}
